%! Introduction to Functions and Change — Unit 1, Lesson 1.0 — Homework (ANSWER KEY)
% THE LESSON'S GRADED PRACTICE and the last component in the packet — exactly TWO pages, started
% in the Homework Launch (the last 3 minutes of the period) and finished at home. Every context
% is new: the notes used a phone plan, a candle, a bike ride, twins, and a runner; the homework
% uses a phone charger, a café, a bathtub, and a parking garage.
%   Page 1 — Part A (input and output, 1–3) and Part B (function or not, 4–7).
%   Page 2 — Part C (a graph, 8–10) and Part D (a table and a flip, 11–12).
% The diagnostic item is 7: "many students chose blue, so it's not a function" — the
% many-to-one misconception, in a new costume. Problem 12 asks for the headline in writing.
%
% PAGE LOCKSTEP with the other half of the pair: prose answers are \writespace{H}{…} (fixed
% height both ways), and each \blank sits at the END of a line, or in a table cell, so its
% \ans cannot rewrap anything. The explicit \newpage fixes the page break in both files.
\documentclass[10pt]{article}
\usepackage{precalculus-article}
\usepackage{precalculus-key}
\newcommand{\TallMath}[1]{$\displaystyle #1\rule[-1.4em]{0pt}{3.2em}$}

\begin{document}

\pageheader{Unit 1, Lesson 1.0}{Homework --- Answer Key}

\vspace{0.12in}

\boxguard[8]
\begin{remindbox}
\small
\textbf{This is your graded homework.} Start it in class; finish it on your own by the due
date on your cover sheet. For every function question, check the \emph{inputs}: does each one
have exactly one output?
\end{remindbox}

\vspace{0.12in}

\begin{headlinebox}{lilac}
\textbf{\color{plum}Part A --- Input and output}
\end{headlinebox}

\begin{enumerate}[leftmargin=1.9em, itemsep=8pt]

\item A phone is plugged in to charge, and its battery percent is checked every few minutes.

      Input: \ans{minutes plugged in} \qquad Output: \ans{battery percent}

\item A café tracks the outside temperature each day and the number of hot chocolates it sells.

      Input: \ans{outside temperature} \qquad Output: \ans{hot chocolates sold}

      Explain your choice using the ``you-control'' idea: which quantity responds to the other?
      \writespace{1.1cm}{The number of hot chocolates sold responds to the temperature; the temperature does not respond to sales.}

\item For the phone in problem 1, the ordered pair $(30, 45)$ is recorded. Write what it means in
      a sentence.
      \writespace{1.1cm}{After 30 minutes plugged in, the battery is at 45 percent.}

\end{enumerate}

\vspace{0.1in}

\begin{headlinebox}{lilac}
\textbf{\color{plum}Part B --- Function or not?}
\end{headlinebox}

\begin{enumerate}[leftmargin=1.9em, itemsep=8pt, start=4]

\item Write \textbf{yes} or \textbf{no}: is the rule a function?

      \smallskip
      \renewcommand{\arraystretch}{1.35}
      \begin{tabular}{|p{0.36\linewidth}|c||p{0.36\linewidth}|c|}
      \hline
      each person $\to$ their age in years & \ans{yes} &
      each age $\to$ a person that age & \ans{no} \\ \hline
      each word $\to$ its number of letters & \ans{yes} &
      each teacher $\to$ a class they teach & \ans{no} \\ \hline
      \end{tabular}

\item \begin{minipage}[t]{0.58\linewidth}
      The arrow diagram shows a rule. Is it a function? Explain by counting arrows.
      \writespace{1.5cm}{Yes. Each input has exactly one arrow leaving it. A and B each receive two arrows, but that is many-to-one --- allowed.}
      \end{minipage}\hfill
      \begin{minipage}[t]{0.36\linewidth}
      \vspace{-4pt}
      \centering
      \begin{tikzpicture}[>=stealth, font=\small, baseline=(current bounding box.north)]
        \foreach \n/\y in {1/0.9, 2/0.3, 3/-0.3, 4/-0.9} \node (i\n) at (0,\y) {\n};
        \node (oA) at (2.2,0.6) {A};
        \node (oB) at (2.2,-0.6) {B};
        \draw[->, thick, plum] (i1) -- (oA);
        \draw[->, thick, plum] (i2) -- (oA);
        \draw[->, thick, plum] (i3) -- (oB);
        \draw[->, thick, plum] (i4) -- (oB);
      \end{tikzpicture}
      \end{minipage}

\item The pairs $(2, 7),\ (5, 1),\ (2, 4)$ are (input, output). Is this a function? Explain.
      \writespace{1.1cm}{No. The input 2 has two outputs, 7 and 4 --- one-to-many.}

\item Jordan says, ``\emph{Each student $\to$ their favorite color} is not a function, because
      lots of students picked blue.'' Is Jordan right? Explain.
      \writespace{1.5cm}{No. Each student still has exactly one favorite color. Many students sharing the output blue is many-to-one, which is allowed.}

\end{enumerate}

\newpage

\begin{headlinebox}{lilac}
\textbf{\color{plum}Part C --- A graph tells a story}
\end{headlinebox}

\begin{minipage}[t]{0.47\linewidth}
\vspace{0pt}
\centering
\begin{tikzpicture}[x=0.42cm, y=0.4cm]
  \draw[linegray, very thin, step=1] (0,0) grid (14,13);
  \draw[->, thick] (0,0) -- (14.6,0) node[below right] {\small min};
  \draw[->, thick] (0,0) -- (0,13.6) node[above] {\small water depth (in)};
  \foreach \x in {2,4,6,8,10,12,14} \node[below, font=\scriptsize] at (\x,0) {\x};
  \foreach \y in {3,6,9,12} \node[left, font=\scriptsize] at (0,\y) {\y};
  \draw[very thick, plum] (0,0) -- (4,12) -- (10,12) -- (14,0);
\end{tikzpicture}

\small A bathtub fills, someone takes a bath, then it drains.
\end{minipage}\hfill
\begin{minipage}[t]{0.5\linewidth}
\vspace{0pt}
\begin{enumerate}[leftmargin=1.9em, itemsep=8pt, start=8]
\item Name each interval.

      Increasing: from \ans{0 to 4 minutes}

      Constant: from \ans{4 to 10 minutes}

      Decreasing: from \ans{10 to 14 minutes}

\item What is the depth at 7 minutes, and what is happening in the tub then?
      \writespace{1.4cm}{12 inches. Someone is taking a bath: time passes, but the depth holds steady.}

\item The depth is 6 inches at 2 minutes and again at 12 minutes. Is depth a function of time?
      Is time a function of depth? Explain both.
      \writespace{2.3cm}{Depth IS a function of time: each time has one depth (two times sharing 6 in is many-to-one). Time is NOT a function of depth: the input 6 in has two outputs, 2 and 12 minutes.}
\end{enumerate}
\end{minipage}

\vspace{0.1in}

\begin{headlinebox}{lilac}
\textbf{\color{plum}Part D --- A table, and the headline}
\end{headlinebox}

\begin{enumerate}[leftmargin=1.9em, itemsep=8pt, start=11]

\item A parking garage charges by the hour, up to a daily maximum.

      \smallskip\centerline{\renewcommand{\arraystretch}{1.3}%
      \begin{tabular}{|c|c|c|c|c|c|}
      \hline
      Hours parked & 1 & 2 & 3  & 4  & 5 \\ \hline
      Cost (\$)    & 5 & 8 & 11 & 11 & 11 \\ \hline
      \end{tabular}}
      \smallskip

      \textbf{(a)} The cost is increasing from hour \ans{1 to hour 3}

      \textbf{(b)} The cost is constant from hour \ans{3 to hour 5}

      \textbf{(c)} A friend paid \$11. Can you tell exactly how many hours they parked? What
      does that say about \emph{hours} as a function of \emph{cost}?
      \writespace{1.6cm}{No --- \$11 could mean 3, 4, or 5 hours. The input \$11 has several outputs, so hours is not a function of cost (even though cost is a function of hours).}

\item In your own words: what makes a rule a function? Then say which one is allowed ---
      two inputs sharing an output, or one input with two outputs.
      \writespace{1.6cm}{Each input has exactly one output. Two inputs sharing an output is allowed; one input with two outputs is not.}

\end{enumerate}

\vspace{0.04in}
\begin{spiralbox}
\small
\textbf{Next lesson (1.1, Function Fundamentals):} the machine gets a name. Instead of writing
``cost is a function of hours parked,'' we will write $C(3) = 11$ --- one short symbol for the
input, the output, and the rule.
\end{spiralbox}

\end{document}
